En un banco de ensayo simulado, una pieza de \(1\,\mathrm{kg}\) se mueve horizontalmente. Si \(x\) es su desplazamiento desde el equilibrio, el resorte ejerce \(-kx\), con \(k=1\,\mathrm{N/m}\), y el amortiguador ejerce \(-cx'\), con \(c=2\,\mathrm{N\,s/m}\). Inicialmente \(x(0)=x'(0)=0\). La fuerza externa es \(1\,\mathrm N\) durante el primer segundo, \(b\,\mathrm N\) durante el segundo y cero después; \(b\) puede tener cualquier signo. No hay golpes ni impulsos ideales.

El accionamiento debe conseguir que la primera parada después de arrancar ocurra en \(t=2\,\mathrm s\). Construya un modelo del movimiento que permita determinar \(b\) y elija un procedimiento de resolución, justificándolo. Calcule el nivel de fuerza necesario y el desplazamiento en ese instante. Determine también cómo evoluciona la pieza después de apagar la fuerza.
